\section{结果与讨论框架}
\textit{本节用于组织正式投稿稿件中的结果呈现。当前版本给出建议的小节结构与需要回答的核心问题，待实验完成后补充图表、误差统计和讨论文本。}

\subsection{接触斑重建质量}
本小节应首先报告接触斑面积、接触中心位置和法向分布的一致性，以回答所提出方法是否能够稳定捕捉面接触区域及其时变演化。对于可获得解析解或高精度参考解的算例，应给出接触斑几何误差随时间的变化。

\subsection{净接触力与净接触力矩预测}
本小节应对比各方法在法向合力、力矩及其峰值时刻上的预测差异，重点说明所提出方法在偏心受压、边角接触和非凸接触中是否能更准确地恢复分布式接触力矩。

\subsection{动态响应连续性与数值稳定性}
本小节应展示接触切换阶段的力时程、速度时程与能量演化，分析高频抖动、非物理跳变和时间步敏感性。对复杂接触切换场景，应特别关注接触激活/失活时响应是否连续。

\subsection{参数敏感性与计算成本}
本小节应报告 SDF 分辨率、表面采样密度、时间步长、接触阈值和稳定化参数对结果的影响，同时给出单步耗时、总耗时以及与基线模型相比的额外计算代价。
